% --- Archon physics formalization source begin ---
% archon:physics
% archon:covers PhyXMiniProblems/problem_phyx_mini_0359.lean
% archon:source-report reports/phyx_mini/problem_phyx_mini_0359.source.json

\chapter{Physics problem phyx\_mini\_0359}
\label{ch:PhyXMiniProblems_problem_phyx_mini_0359}

\paragraph{Problem source.}
\#\# Physical scenario

A cylindrical container 80 cm long is separated into two compartments by a thin piston, originally clamped in position 30 cm from the left end. The left compartment is filled with one mole of helium gas at a pressure of 5 atmospheres; the right compartment is filled with argon gas at 1 atmosphere of pressure. These gases may be considered ideal. The cylinder is submerged in 1 liter of water, and the entire system is initially at the uniform temperature of \textbackslash{}\textasciicircum{}\textbackslash{}circ C. The heat capacities of the cylinder and piston may be neglected. When the piston is unclamped, a new equilibrium situation is ultimately reached with the piston in a new position.

\#\# Question

How far from the left end of the cylinder will the piston come to rest?

\#\# Answer choices

- A: A: 60cm

- B: B: 30cm

- C: C: 50cm

- D: D: 40cm

\#\# Figure caption (auxiliary; use the image as primary evidence)

The image depicts a cross-sectional diagram of a container submerged in water. The container is divided into two compartments by a vertical partition.

- **Left Compartment:** Labeled "He" with a pressure of "5 atm". It measures 30 cm in width.
- **Right Compartment:** Labeled "Ar" with a pressure of "1 atm". It measures 50 cm in width.

The entire container appears to be surrounded by water, as indicated by the label "Water" on the right side.

\#\# Dataset metadata

Ideal Gases and Kinetic Theory; Physical Model Grounding Reasoning, Multi-Formula Reasoning

\paragraph{Recorded answer/context.}
A: A: 60cm

\paragraph{Figure/image path.}
/root/proposal\_for\_physic/science-mango/phyx\_mini\_run/phyx\_data/test\_image/359.png

\paragraph{Formalization target.}
create a compiling Lean file with sorry bodies at `PhyXMiniProblems/problem\_phyx\_mini\_0359.lean`. The Lean declarations must preserve the physical quantities, dimensions or dimensional roles, figure labels, governing-law hypotheses, and final relation expressed by this problem.
Use Mathlib/Physlib names found through LeanExplore where available. If a physics API is missing, introduce faithful local abstractions rather than scalar placeholder aliases.

\begin{theorem}[Physics formalization target]
\label{thm:physics:phyx_mini_0359:target}
\lean{PhyXMiniProblems.ProblemPhyXMini0359.problem_phyx_mini_0359}
\uses{def:physics:phyx-mini-0359:phyxminiproblems-problemphyxmini0359-lengthincentimeters, def:physics:phyx-mini-0359:phyxminiproblems-problemphyxmini0359-twogaspistonsetup, def:physics:phyx-mini-0359:phyxminiproblems-problemphyxmini0359-matchesproblemscenario, def:physics:phyx-mini-0359:phyxminiproblems-problemphyxmini0359-matchesproblemandfiguredata, def:physics:phyx-mini-0359:phyxminiproblems-problemphyxmini0359-hasphysicalparameters, def:physics:phyx-mini-0359:phyxminiproblems-problemphyxmini0359-satisfiesidealgasandequilibriumlaws, def:physics:phyx-mini-0359:phyxminiproblems-problemphyxmini0359-isuniquematchinganswer}
The assigned autoformalize agent should translate this physics problem into Lean declarations in the covered file, with theorem and lemma proof bodies written as `by sorry`.
\end{theorem}
\begin{proof}
This is an autoformalization task, not a proof task. Produce statements that can later be proved by the physics prover without weakening the physical model.
\end{proof}

% --- Archon declaration topology begin ---
\section{Declaration topology}
\begin{definition}[Length Quantity]
  \label{def:physics:phyx-mini-0359:phyxminiproblems-problemphyxmini0359-lengthquantity}
  \lean{PhyXMiniProblems.ProblemPhyXMini0359.LengthQuantity}
  A nonnegative physical length, independent of the unit used to read it
\end{definition}
\begin{proof}
  This is the declared model definition of Length Quantity.
\end{proof}
\begin{definition}[Area Quantity]
  \label{def:physics:phyx-mini-0359:phyxminiproblems-problemphyxmini0359-areaquantity}
  \lean{PhyXMiniProblems.ProblemPhyXMini0359.AreaQuantity}
  A nonnegative physical cross-sectional area
\end{definition}
\begin{proof}
  This is the declared model definition of Area Quantity.
\end{proof}
\begin{definition}[Volume Quantity]
  \label{def:physics:phyx-mini-0359:phyxminiproblems-problemphyxmini0359-volumequantity}
  \lean{PhyXMiniProblems.ProblemPhyXMini0359.VolumeQuantity}
  A nonnegative physical volume with dimension length cubed
\end{definition}
\begin{proof}
  This is the declared model definition of Volume Quantity.
\end{proof}
\begin{definition}[length Readout]
  \label{def:physics:phyx-mini-0359:phyxminiproblems-problemphyxmini0359-lengthreadout}
  \lean{PhyXMiniProblems.ProblemPhyXMini0359.lengthReadout}
  \uses{def:physics:phyx-mini-0359:phyxminiproblems-problemphyxmini0359-lengthquantity}
  Read a physical length in a selected length unit
\end{definition}
\begin{proof}
  This is the declared model definition of length Readout.
\end{proof}
\begin{definition}[area Readout]
  \label{def:physics:phyx-mini-0359:phyxminiproblems-problemphyxmini0359-areareadout}
  \lean{PhyXMiniProblems.ProblemPhyXMini0359.areaReadout}
  \uses{def:physics:phyx-mini-0359:phyxminiproblems-problemphyxmini0359-areaquantity}
  Read an area in the square of a selected length unit
\end{definition}
\begin{proof}
  This is the declared model definition of area Readout.
\end{proof}
\begin{definition}[volume Readout]
  \label{def:physics:phyx-mini-0359:phyxminiproblems-problemphyxmini0359-volumereadout}
  \lean{PhyXMiniProblems.ProblemPhyXMini0359.volumeReadout}
  \uses{def:physics:phyx-mini-0359:phyxminiproblems-problemphyxmini0359-volumequantity}
  Read a volume in the cube of a selected length unit
\end{definition}
\begin{proof}
  This is the declared model definition of volume Readout.
\end{proof}
\begin{definition}[pressure Readout]
  \label{def:physics:phyx-mini-0359:phyxminiproblems-problemphyxmini0359-pressurereadout}
  \lean{PhyXMiniProblems.ProblemPhyXMini0359.pressureReadout}
  Read a pressure in coherent units selected by units
\end{definition}
\begin{proof}
  This is the declared model definition of pressure Readout.
\end{proof}
\begin{definition}[length In Centimeters]
  \label{def:physics:phyx-mini-0359:phyxminiproblems-problemphyxmini0359-lengthincentimeters}
  \lean{PhyXMiniProblems.ProblemPhyXMini0359.lengthInCentimeters}
  \uses{def:physics:phyx-mini-0359:phyxminiproblems-problemphyxmini0359-lengthquantity, def:physics:phyx-mini-0359:phyxminiproblems-problemphyxmini0359-lengthreadout}
  Centimeter readout used for the cylinder and piston coordinates
\end{definition}
\begin{proof}
  This is the declared model definition of length In Centimeters.
\end{proof}
\begin{definition}[area In Square Centimeters]
  \label{def:physics:phyx-mini-0359:phyxminiproblems-problemphyxmini0359-areainsquarecentimeters}
  \lean{PhyXMiniProblems.ProblemPhyXMini0359.areaInSquareCentimeters}
  \uses{def:physics:phyx-mini-0359:phyxminiproblems-problemphyxmini0359-areaquantity, def:physics:phyx-mini-0359:phyxminiproblems-problemphyxmini0359-areareadout}
  Square-centimeter readout of the cylinder cross-sectional area
\end{definition}
\begin{proof}
  This is the declared model definition of area In Square Centimeters.
\end{proof}
\begin{definition}[volume In Cubic Centimeters]
  \label{def:physics:phyx-mini-0359:phyxminiproblems-problemphyxmini0359-volumeincubiccentimeters}
  \lean{PhyXMiniProblems.ProblemPhyXMini0359.volumeInCubicCentimeters}
  \uses{def:physics:phyx-mini-0359:phyxminiproblems-problemphyxmini0359-volumequantity, def:physics:phyx-mini-0359:phyxminiproblems-problemphyxmini0359-volumereadout}
  Cubic-centimeter readout of a physical volume
\end{definition}
\begin{proof}
  This is the declared model definition of volume In Cubic Centimeters.
\end{proof}
\begin{definition}[volume In Liters]
  \label{def:physics:phyx-mini-0359:phyxminiproblems-problemphyxmini0359-volumeinliters}
  \lean{PhyXMiniProblems.ProblemPhyXMini0359.volumeInLiters}
  \uses{def:physics:phyx-mini-0359:phyxminiproblems-problemphyxmini0359-volumequantity, def:physics:phyx-mini-0359:phyxminiproblems-problemphyxmini0359-volumeincubiccentimeters}
  Liter readout, using 1 L = 1000 cm³
\end{definition}
\begin{proof}
  This is the declared model definition of volume In Liters.
\end{proof}
\begin{definition}[pressure In Atmospheres]
  \label{def:physics:phyx-mini-0359:phyxminiproblems-problemphyxmini0359-pressureinatmospheres}
  \lean{PhyXMiniProblems.ProblemPhyXMini0359.pressureInAtmospheres}
  \uses{def:physics:phyx-mini-0359:phyxminiproblems-problemphyxmini0359-pressurereadout}
  Atmosphere readout formed from Physlib's dimensionful standard atmosphere, rather than identifying pressure with a bare real number
\end{definition}
\begin{proof}
  This is the declared model definition of pressure In Atmospheres.
\end{proof}
\begin{definition}[temperature In Kelvin]
  \label{def:physics:phyx-mini-0359:phyxminiproblems-problemphyxmini0359-temperatureinkelvin}
  \lean{PhyXMiniProblems.ProblemPhyXMini0359.temperatureInKelvin}
  Read an absolute Temperature in kelvin. Physlib stores a temperature in an arbitrary zero-preserving unit, so the storage unit remains explicit
\end{definition}
\begin{proof}
  This is the declared model definition of temperature In Kelvin.
\end{proof}
\begin{definition}[kelvin To Celsius]
  \label{def:physics:phyx-mini-0359:phyxminiproblems-problemphyxmini0359-kelvintocelsius}
  \lean{PhyXMiniProblems.ProblemPhyXMini0359.kelvinToCelsius}
  Convert an absolute kelvin readout to degrees Celsius
\end{definition}
\begin{proof}
  This is the declared model definition of kelvin To Celsius.
\end{proof}
\begin{definition}[Gas Species]
  \label{def:physics:phyx-mini-0359:phyxminiproblems-problemphyxmini0359-gasspecies}
  \lean{PhyXMiniProblems.ProblemPhyXMini0359.GasSpecies}
  The two gases named inside the compartments
\end{definition}
\begin{proof}
  This is the declared model definition of Gas Species.
\end{proof}
\begin{definition}[System State]
  \label{def:physics:phyx-mini-0359:phyxminiproblems-problemphyxmini0359-systemstate}
  \lean{PhyXMiniProblems.ProblemPhyXMini0359.SystemState}
  Initial clamped state and final equilibrium state
\end{definition}
\begin{proof}
  This is the declared model definition of System State.
\end{proof}
\begin{definition}[Thermal Body]
  \label{def:physics:phyx-mini-0359:phyxminiproblems-problemphyxmini0359-thermalbody}
  \lean{PhyXMiniProblems.ProblemPhyXMini0359.ThermalBody}
  Bodies whose temperatures are relevant to thermal equilibrium
\end{definition}
\begin{proof}
  This is the declared model definition of Thermal Body.
\end{proof}
\begin{definition}[gas Thermal Body]
  \label{def:physics:phyx-mini-0359:phyxminiproblems-problemphyxmini0359-gasthermalbody}
  \lean{PhyXMiniProblems.ProblemPhyXMini0359.gasThermalBody}
  \uses{def:physics:phyx-mini-0359:phyxminiproblems-problemphyxmini0359-gasspecies, def:physics:phyx-mini-0359:phyxminiproblems-problemphyxmini0359-thermalbody}
  Associate each gas sample with its thermal body
\end{definition}
\begin{proof}
  This is the declared model definition of gas Thermal Body.
\end{proof}
\begin{definition}[Gas Model]
  \label{def:physics:phyx-mini-0359:phyxminiproblems-problemphyxmini0359-gasmodel}
  \lean{PhyXMiniProblems.ProblemPhyXMini0359.GasModel}
  The ideal/nonideal gas-model distinction stated in the prose
\end{definition}
\begin{proof}
  This is the declared model definition of Gas Model.
\end{proof}
\begin{definition}[Container Model]
  \label{def:physics:phyx-mini-0359:phyxminiproblems-problemphyxmini0359-containermodel}
  \lean{PhyXMiniProblems.ProblemPhyXMini0359.ContainerModel}
  Geometrical idealization of the vessel
\end{definition}
\begin{proof}
  This is the declared model definition of Container Model.
\end{proof}
\begin{definition}[Piston Model]
  \label{def:physics:phyx-mini-0359:phyxminiproblems-problemphyxmini0359-pistonmodel}
  \lean{PhyXMiniProblems.ProblemPhyXMini0359.PistonModel}
  Mechanical idealization of the separating piston
\end{definition}
\begin{proof}
  This is the declared model definition of Piston Model.
\end{proof}
\begin{definition}[Piston Constraint]
  \label{def:physics:phyx-mini-0359:phyxminiproblems-problemphyxmini0359-pistonconstraint}
  \lean{PhyXMiniProblems.ProblemPhyXMini0359.PistonConstraint}
  Constraint imposed on the piston at a stage of the experiment
\end{definition}
\begin{proof}
  This is the declared model definition of Piston Constraint.
\end{proof}
\begin{definition}[Heat Capacity Treatment]
  \label{def:physics:phyx-mini-0359:phyxminiproblems-problemphyxmini0359-heatcapacitytreatment}
  \lean{PhyXMiniProblems.ProblemPhyXMini0359.HeatCapacityTreatment}
  Whether a solid body's heat capacity is retained in the thermal model
\end{definition}
\begin{proof}
  This is the declared model definition of Heat Capacity Treatment.
\end{proof}
\begin{definition}[Piston Cylinder Figure]
  \label{def:physics:phyx-mini-0359:phyxminiproblems-problemphyxmini0359-pistoncylinderfigure}
  \lean{PhyXMiniProblems.ProblemPhyXMini0359.PistonCylinderFigure}
  \uses{def:physics:phyx-mini-0359:phyxminiproblems-problemphyxmini0359-lengthquantity, def:physics:phyx-mini-0359:phyxminiproblems-problemphyxmini0359-gasspecies}
  Information read from the primary image. It distinguishes the left and right regions, their gas labels, displayed pressures, widths, and the water surrounding the cylinder. It contains no final piston position
\end{definition}
\begin{proof}
  This is the declared model definition of Piston Cylinder Figure.
\end{proof}
\begin{definition}[Two Gas Piston Setup]
  \label{def:physics:phyx-mini-0359:phyxminiproblems-problemphyxmini0359-twogaspistonsetup}
  \lean{PhyXMiniProblems.ProblemPhyXMini0359.TwoGasPistonSetup}
  \uses{def:physics:phyx-mini-0359:phyxminiproblems-problemphyxmini0359-lengthquantity, def:physics:phyx-mini-0359:phyxminiproblems-problemphyxmini0359-areaquantity, def:physics:phyx-mini-0359:phyxminiproblems-problemphyxmini0359-volumequantity, def:physics:phyx-mini-0359:phyxminiproblems-problemphyxmini0359-gasspecies, def:physics:phyx-mini-0359:phyxminiproblems-problemphyxmini0359-systemstate, def:physics:phyx-mini-0359:phyxminiproblems-problemphyxmini0359-thermalbody, def:physics:phyx-mini-0359:phyxminiproblems-problemphyxmini0359-gasmodel, def:physics:phyx-mini-0359:phyxminiproblems-problemphyxmini0359-containermodel, def:physics:phyx-mini-0359:phyxminiproblems-problemphyxmini0359-pistonmodel, def:physics:phyx-mini-0359:phyxminiproblems-problemphyxmini0359-pistonconstraint, def:physics:phyx-mini-0359:phyxminiproblems-problemphyxmini0359-heatcapacitytreatment, def:physics:phyx-mini-0359:phyxminiproblems-problemphyxmini0359-pistoncylinderfigure}
  Independent physical quantities and observables of the experiment. Physlib currently has no amount-of-substance dimension, so the amount type and its calibrated mole readout are explicit parameters. A single sealed sample of each gas is used at both states. The final piston position is an independent physical observable, not a definition involving an answer choice
\end{definition}
\begin{proof}
  This is the declared model definition of Two Gas Piston Setup.
\end{proof}
\begin{definition}[gas Amount In Moles]
  \label{def:physics:phyx-mini-0359:phyxminiproblems-problemphyxmini0359-gasamountinmoles}
  \lean{PhyXMiniProblems.ProblemPhyXMini0359.gasAmountInMoles}
  \uses{def:physics:phyx-mini-0359:phyxminiproblems-problemphyxmini0359-gasspecies, def:physics:phyx-mini-0359:phyxminiproblems-problemphyxmini0359-twogaspistonsetup}
  Mole readout of one of the two sealed gas samples
\end{definition}
\begin{proof}
  This is the declared model definition of gas Amount In Moles.
\end{proof}
\begin{definition}[body Temperature In Kelvin]
  \label{def:physics:phyx-mini-0359:phyxminiproblems-problemphyxmini0359-bodytemperatureinkelvin}
  \lean{PhyXMiniProblems.ProblemPhyXMini0359.bodyTemperatureInKelvin}
  \uses{def:physics:phyx-mini-0359:phyxminiproblems-problemphyxmini0359-temperatureinkelvin, def:physics:phyx-mini-0359:phyxminiproblems-problemphyxmini0359-systemstate, def:physics:phyx-mini-0359:phyxminiproblems-problemphyxmini0359-thermalbody, def:physics:phyx-mini-0359:phyxminiproblems-problemphyxmini0359-twogaspistonsetup}
  Kelvin readout of a body's temperature in one system state
\end{definition}
\begin{proof}
  This is the declared model definition of body Temperature In Kelvin.
\end{proof}
\begin{definition}[body Temperature In Celsius]
  \label{def:physics:phyx-mini-0359:phyxminiproblems-problemphyxmini0359-bodytemperatureincelsius}
  \lean{PhyXMiniProblems.ProblemPhyXMini0359.bodyTemperatureInCelsius}
  \uses{def:physics:phyx-mini-0359:phyxminiproblems-problemphyxmini0359-kelvintocelsius, def:physics:phyx-mini-0359:phyxminiproblems-problemphyxmini0359-systemstate, def:physics:phyx-mini-0359:phyxminiproblems-problemphyxmini0359-thermalbody, def:physics:phyx-mini-0359:phyxminiproblems-problemphyxmini0359-twogaspistonsetup, def:physics:phyx-mini-0359:phyxminiproblems-problemphyxmini0359-bodytemperatureinkelvin}
  Celsius readout of a body's temperature in one system state
\end{definition}
\begin{proof}
  This is the declared model definition of body Temperature In Celsius.
\end{proof}
\begin{definition}[chamber Width In Centimeters]
  \label{def:physics:phyx-mini-0359:phyxminiproblems-problemphyxmini0359-chamberwidthincentimeters}
  \lean{PhyXMiniProblems.ProblemPhyXMini0359.chamberWidthInCentimeters}
  \uses{def:physics:phyx-mini-0359:phyxminiproblems-problemphyxmini0359-lengthincentimeters, def:physics:phyx-mini-0359:phyxminiproblems-problemphyxmini0359-gasspecies, def:physics:phyx-mini-0359:phyxminiproblems-problemphyxmini0359-systemstate, def:physics:phyx-mini-0359:phyxminiproblems-problemphyxmini0359-twogaspistonsetup}
  Axial width of a gas compartment. Helium occupies the portion to the left of the piston and argon the remaining portion of the rigid cylinder
\end{definition}
\begin{proof}
  This is the declared model definition of chamber Width In Centimeters.
\end{proof}
\begin{definition}[chamber Volume In Cubic Centimeters]
  \label{def:physics:phyx-mini-0359:phyxminiproblems-problemphyxmini0359-chambervolumeincubiccentimeters}
  \lean{PhyXMiniProblems.ProblemPhyXMini0359.chamberVolumeInCubicCentimeters}
  \uses{def:physics:phyx-mini-0359:phyxminiproblems-problemphyxmini0359-areainsquarecentimeters, def:physics:phyx-mini-0359:phyxminiproblems-problemphyxmini0359-gasspecies, def:physics:phyx-mini-0359:phyxminiproblems-problemphyxmini0359-systemstate, def:physics:phyx-mini-0359:phyxminiproblems-problemphyxmini0359-twogaspistonsetup, def:physics:phyx-mini-0359:phyxminiproblems-problemphyxmini0359-chamberwidthincentimeters}
  Cylindrical chamber volume area * axial width, read coherently in cubic centimeters. This is the rigid-cylinder geometry law, not a condition on the requested equilibrium position
\end{definition}
\begin{proof}
  This is the declared model definition of chamber Volume In Cubic Centimeters.
\end{proof}
\begin{definition}[Obeys Ideal Gas Law]
  \label{def:physics:phyx-mini-0359:phyxminiproblems-problemphyxmini0359-obeysidealgaslaw}
  \lean{PhyXMiniProblems.ProblemPhyXMini0359.ObeysIdealGasLaw}
  The ideal-gas equation in one coherent collection of named-unit readouts. The gas constant therefore has units atm cm³ mol⁻¹ K⁻¹
\end{definition}
\begin{proof}
  This is the declared model definition of Obeys Ideal Gas Law.
\end{proof}
\begin{definition}[Matches Problem Scenario]
  \label{def:physics:phyx-mini-0359:phyxminiproblems-problemphyxmini0359-matchesproblemscenario}
  \lean{PhyXMiniProblems.ProblemPhyXMini0359.MatchesProblemScenario}
  \uses{def:physics:phyx-mini-0359:phyxminiproblems-problemphyxmini0359-twogaspistonsetup}
  Qualitative idealizations and experimental sequence stated in the prose
\end{definition}
\begin{proof}
  This is the declared model definition of Matches Problem Scenario.
\end{proof}
\begin{definition}[Matches Problem And Figure Data]
  \label{def:physics:phyx-mini-0359:phyxminiproblems-problemphyxmini0359-matchesproblemandfiguredata}
  \lean{PhyXMiniProblems.ProblemPhyXMini0359.MatchesProblemAndFigureData}
  \uses{def:physics:phyx-mini-0359:phyxminiproblems-problemphyxmini0359-lengthincentimeters, def:physics:phyx-mini-0359:phyxminiproblems-problemphyxmini0359-volumeinliters, def:physics:phyx-mini-0359:phyxminiproblems-problemphyxmini0359-pressureinatmospheres, def:physics:phyx-mini-0359:phyxminiproblems-problemphyxmini0359-twogaspistonsetup, def:physics:phyx-mini-0359:phyxminiproblems-problemphyxmini0359-gasamountinmoles, def:physics:phyx-mini-0359:phyxminiproblems-problemphyxmini0359-bodytemperatureincelsius}
  Numerical problem data and primary-image readouts. The source report restores the garbled temperature in the raw text as 25 °C. No final coordinate or answer choice occurs in these fields
\end{definition}
\begin{proof}
  This is the declared model definition of Matches Problem And Figure Data.
\end{proof}
\begin{definition}[Has Physical Parameters]
  \label{def:physics:phyx-mini-0359:phyxminiproblems-problemphyxmini0359-hasphysicalparameters}
  \lean{PhyXMiniProblems.ProblemPhyXMini0359.HasPhysicalParameters}
  \uses{def:physics:phyx-mini-0359:phyxminiproblems-problemphyxmini0359-lengthincentimeters, def:physics:phyx-mini-0359:phyxminiproblems-problemphyxmini0359-areainsquarecentimeters, def:physics:phyx-mini-0359:phyxminiproblems-problemphyxmini0359-pressureinatmospheres, def:physics:phyx-mini-0359:phyxminiproblems-problemphyxmini0359-twogaspistonsetup, def:physics:phyx-mini-0359:phyxminiproblems-problemphyxmini0359-gasamountinmoles, def:physics:phyx-mini-0359:phyxminiproblems-problemphyxmini0359-bodytemperatureinkelvin}
  Positivity and nondegeneracy conditions selecting physical gas states
\end{definition}
\begin{proof}
  This is the declared model definition of Has Physical Parameters.
\end{proof}
\begin{definition}[Satisfies Ideal Gas And Equilibrium Laws]
  \label{def:physics:phyx-mini-0359:phyxminiproblems-problemphyxmini0359-satisfiesidealgasandequilibriumlaws}
  \lean{PhyXMiniProblems.ProblemPhyXMini0359.SatisfiesIdealGasAndEquilibriumLaws}
  \uses{def:physics:phyx-mini-0359:phyxminiproblems-problemphyxmini0359-pressureinatmospheres, def:physics:phyx-mini-0359:phyxminiproblems-problemphyxmini0359-gasthermalbody, def:physics:phyx-mini-0359:phyxminiproblems-problemphyxmini0359-twogaspistonsetup, def:physics:phyx-mini-0359:phyxminiproblems-problemphyxmini0359-gasamountinmoles, def:physics:phyx-mini-0359:phyxminiproblems-problemphyxmini0359-bodytemperatureinkelvin, def:physics:phyx-mini-0359:phyxminiproblems-problemphyxmini0359-chambervolumeincubiccentimeters, def:physics:phyx-mini-0359:phyxminiproblems-problemphyxmini0359-obeysidealgaslaw}
  Governing laws for the sealed ideal gases and their final equilibrium. both gas samples obey PV = nRT before and after release; the same gasSample gas occurs at both states, expressing amount conservation in each sealed compartment; final mechanical equilibrium gives equal gas pressures; final thermal equilibrium gives one common temperature for helium, argon, and the surrounding water. These laws are generic and contain no numerical final piston position
\end{definition}
\begin{proof}
  This is the declared model definition of Satisfies Ideal Gas And Equilibrium Laws.
\end{proof}
\begin{definition}[Answer Choice]
  \label{def:physics:phyx-mini-0359:phyxminiproblems-problemphyxmini0359-answerchoice}
  \lean{PhyXMiniProblems.ProblemPhyXMini0359.AnswerChoice}
  Labels printed beside the four proposed piston positions
\end{definition}
\begin{proof}
  This is the declared model definition of Answer Choice.
\end{proof}
\begin{definition}[displayed Piston Position In Centimeters]
  \label{def:physics:phyx-mini-0359:phyxminiproblems-problemphyxmini0359-displayedpistonpositionincentimeters}
  \lean{PhyXMiniProblems.ProblemPhyXMini0359.displayedPistonPositionInCentimeters}
  \uses{def:physics:phyx-mini-0359:phyxminiproblems-problemphyxmini0359-answerchoice}
  Piston distance from the left end, in centimeters, printed by each choice
\end{definition}
\begin{proof}
  This is the declared model definition of displayed Piston Position In Centimeters.
\end{proof}
\begin{definition}[Matches Answer Choice]
  \label{def:physics:phyx-mini-0359:phyxminiproblems-problemphyxmini0359-matchesanswerchoice}
  \lean{PhyXMiniProblems.ProblemPhyXMini0359.MatchesAnswerChoice}
  \uses{def:physics:phyx-mini-0359:phyxminiproblems-problemphyxmini0359-answerchoice, def:physics:phyx-mini-0359:phyxminiproblems-problemphyxmini0359-displayedpistonpositionincentimeters}
  A choice exactly matches a derived piston-position readout
\end{definition}
\begin{proof}
  This is the declared model definition of Matches Answer Choice.
\end{proof}
\begin{definition}[Is Unique Matching Answer]
  \label{def:physics:phyx-mini-0359:phyxminiproblems-problemphyxmini0359-isuniquematchinganswer}
  \lean{PhyXMiniProblems.ProblemPhyXMini0359.IsUniqueMatchingAnswer}
  \uses{def:physics:phyx-mini-0359:phyxminiproblems-problemphyxmini0359-answerchoice, def:physics:phyx-mini-0359:phyxminiproblems-problemphyxmini0359-matchesanswerchoice}
  A selected answer is the unique displayed value matching the position
\end{definition}
\begin{proof}
  This is the declared model definition of Is Unique Matching Answer.
\end{proof}
% --- Archon declaration topology end ---
% --- Archon physics formalization source end ---
